\section*{Appendix A. Uniform collaborative research protocol}
\label{app:unified}
\setcounter{table}{0}
\renewcommand{\thetable}{A.\arabic{table}}
\renewcommand{\theHtable}{A.\arabic{table}}
\makeatletter
\setlength{\@fptop}{0pt}
\makeatother

\paragraph{Versioned implementation.}
The reported study uses \texttt{core\_v3.py}, \texttt{run\_v3.py} and \texttt{research\_v2.py}. All task families execute the same optimization, aggregation, proposal and retention source through task-specific adapters. TAPB workers use the TAPB environment; Qwen, proteomics and cell adapters use the LDA environment. Each scientific fit starts from its task-specific public or random initialization, with no reuse of earlier task-trained weights.

\paragraph{Shared optimization and checkpoint selection.}
Each round uses one local AdamW epoch, batch size 64 and gradient-norm clipping at one; local optimizer state resets between rounds. Parameters are averaged with training-example weights. If server momentum is enabled, the coordinator multiplies the previous velocity by the configured momentum coefficient and adds the difference between the averaged parameters and the previous shared parameters, then adds that velocity to the shared parameters. Floating model buffers are averaged and integer buffers use an elementwise maximum. Momentum acts on trainable parameters only. Every full-training candidate runs for 100 rounds; local development scores are collected every five rounds. Checkpoint selection maximizes the equal-client primary score, then minimizes mean development loss. The optional residual adapter has width 32 and zero-initialized output. The DTI adapter trains TAPB using frozen public protein features, the proteomic adapter trains its efficacy head from random initialization, and the cell adapters train response heads on frozen corrected scDEBART features.

\paragraph{Complete design library.}
Each identifier below compiles to one executable candidate. D00 is the fixed starting recipe; D01--D11 are available to both direct and loop search without within-arm repetition. Six proposal slots define the common budget.

\begin{center}\small
\begin{tabular}{lrrrr}
\toprule
Design & Learning rate & Weight decay & Momentum & Residual adapter \\
\midrule
D00 & 0.001 & 0.001 & 0 & no \\
D01 & 0.0001 & 0.001 & 0 & no \\
D02 & 0.01 & 0.001 & 0 & no \\
D03 & 0.001 & 0 & 0 & no \\
D04 & 0.001 & 0.01 & 0 & no \\
D05 & 0.001 & 0.001 & 0.9 & no \\
D06 & 0.001 & 0.001 & 0 & yes \\
D07 & 0.001 & 0.001 & 0.9 & yes \\
D08 & 0.0001 & 0.001 & 0 & yes \\
D09 & 0.0001 & 0.001 & 0.9 & yes \\
D10 & 0.01 & 0.01 & 0 & no \\
D11 & 0.001 & 0.01 & 0 & yes \\
\bottomrule
\end{tabular}
\end{center}

\paragraph{Proposal generation and evidence memory.}
The main-study proposer is Qwen2.5-7B-Instruct, using bfloat16 inference, temperature 0.7, top-$p$ 0.9 and at most 512 new tokens. Each slot uses a deterministic seed derived from the training seed. Direct and loop comparisons share the same first proposal. Later direct proposals receive the default reference configuration and remaining untried identifiers, without evaluation scores. Loop proposals instead receive the latest retained configuration and ordered evidence cards, each containing the candidate configuration, aggregate best development metric and loss, best round, first/last training and validation records, and final status. A bounded second generation attempt repairs schema or novelty violations. Failed proposals consume a slot. Identical configurations within a seed and participation setting can reuse a completed fit while retaining their logical positions in the research traces.

Every proposal has the following executable JSON shape:
\begin{verbatim}
{
  "hypothesis": "one testable claim about an editable factor",
  "experiment": "the controlled change and incumbent comparison",
  "expected_effect": "the expected metric and diagnostic response",
  "design_id": "D00--D11"
}
\end{verbatim}
The compiler validates \texttt{design\_id} against the table above. The finalizer then attaches executed measurements and a retention decision to the proposal, creating the evidence card used in the next iteration.

\paragraph{Executed loop trajectories.}
Three ten-laboratory traces illustrate how the common finalizer operates. Values are aggregate development measurements on the 0--1 primary-metric scale, and proposal slots are one-indexed.

\begin{center}\small
\begin{tabular}{llrrl}
\toprule
Study & Selected design & Primary & Loss & First direct / loop slot \\
\midrule
PTPC, seed 42 & D04 & 0.665833 & 0.521161 & 5 / 5 \\
VCC, seed 43 & D00 & 0.240350 & 1.607792 & initial baseline \\
Norman, seed 42 & D08 & 0.519000 & 1.538505 & 3 / 2 \\
\bottomrule
\end{tabular}
\end{center}

PTPC selects stronger weight decay through the BCE tie-break at equal AP. VCC retains the initial design after six alternatives. On Norman, the loop reaches the winning lower-learning-rate residual design one slot before direct search. Across all 15 ten-laboratory task--seed trajectories, the two-proposal loop is better in two, tied in thirteen and worse in none; the mean development-primary difference is +0.78 percentage points. At six proposals, direct and loop select the same checkpoints, reflecting coverage saturation in the twelve-design library.

\paragraph{Compute-matched evidence racing.}
The follow-up feedback-isolation protocol uses a ten-design fractional-factorial slate and newly executed seeds 53--55. Direct search assigns 80 rounds to every design. The loop assigns 20 screening rounds to every design, stratifies candidates by learning rate, and promotes both the strongest current candidate and the fastest-improving candidate within each stratum. Six promoted candidates restart from the common initialization for 100-round validation. Consequently, each arm executes 800 full-client rounds and 160 development measurements over the same ten hypotheses. The promotion record is finalized before any held-out response is read.

Across the resulting 12 proteomics/cell held-out comparisons, racing gives four wins, eight ties and no losses. Norman improves by 10.84 Top-1 points on average, with seed-level effects of +10.00, +11.70 and +10.81 and a hierarchical paired-bootstrap 95\% interval of [4.54, 17.73]. Tahoe gives one +5.97-point seed and two ties; VCC and PTPC tie in all three seeds. Independent rescoring of every saved prediction reproduces the reported metrics exactly. A TAPB companion replay applies the same screen/promote principle to the seven-design intersection of existing traces: direct receives 525 candidate-rounds, the loop receives 520, and the development AUROC effect is +0.39 points on average (two wins, one tie).

\paragraph{Client and split integrity.}
DTI clients contain 4,932--5,058 training and 517--643 development pairs. Compound identities do not cross clients. The fixed test roles contain 5,505 random, 3,791 unseen-drug and 2,507 unseen-protein pairs. PTPC partitions 333 compound groups into 386 fitting and 105 development conditions, with 35--42 fitting conditions per client; its frozen test role contains 148 conditions. VCC, Norman and Tahoe use whole-intervention splits of 60/20/20, 70/20/27 and 180/60/67. Their ten clients own equal numbers of training and development intervention identities.

\paragraph{Proteomic preprocessing.}
For each 6h and 24h profile, the PTPC adapter applies log1p, within-profile min--max normalization over finite observations, zero filling and unit L2 normalization. The 2,048-bit Morgan vector is separately L2-normalized and scaled by 0.5. Concatenating the two 5,171-dimensional proteomic blocks with the drug vector produces a 12,390-dimensional float32 input. The ProteinTalks-derived efficacy head contains 464,193 parameters at D00. Every arm uses identical inputs, initialization policy, 100-round budget and checkpoint rule.

\paragraph{Full factorial results.}
The following tables report all six held-out arms, complete method contrasts, proposal-prefix analysis, secondary metrics, uncertainty and all three DTI test roles. Bold marks the best displayed value within each task and metric, including ties. Higher is better except BCE and cross-entropy.

\begin{table}[t]
\centering\small
\setlength{\tabcolsep}{3pt}
\begin{tabular}{clrrrrr}
\toprule
 & & DTI & Proteomics & \multicolumn{3}{c}{Cell perturbation} \\
\cmidrule(lr){5-7}
Labs & Method & AUROC & AP & VCC & Norman & Tahoe \\
 & & $\uparrow$ & $\uparrow$ & Top-1 $\uparrow$ & Top-1 $\uparrow$ & Top-1 $\uparrow$ \\
\midrule
1 & Fixed model & 81.76 & 25.05{\scriptsize$\pm4.56$} & 29.94{\scriptsize$\pm0.00$} & \textbf{25.56}{\scriptsize$\pm0.00$} & 19.40{\scriptsize$\pm0.00$} \\
1 & Direct & 88.82 & 20.13{\scriptsize$\pm0.07$} & 21.09{\scriptsize$\pm4.38$} & 14.15{\scriptsize$\pm0.00$} & 15.42{\scriptsize$\pm6.21$} \\
1 & Our loop & 88.82 & 20.13{\scriptsize$\pm0.07$} & 20.22{\scriptsize$\pm2.22$} & 14.15{\scriptsize$\pm0.00$} & 15.42{\scriptsize$\pm6.21$} \\
\midrule
10 & Fixed model & 94.27 & \textbf{37.07}{\scriptsize$\pm0.08$} & \textbf{31.54}{\scriptsize$\pm0.08$} & 15.48{\scriptsize$\pm0.00$} & \textbf{25.87}{\scriptsize$\pm1.72$} \\
10 & Direct & \textbf{94.60} & 36.99{\scriptsize$\pm0.22$} & \textbf{31.54}{\scriptsize$\pm0.08$} & 22.25{\scriptsize$\pm0.45$} & 22.39{\scriptsize$\pm2.59$} \\
10 & Our loop & \textbf{94.60} & 36.99{\scriptsize$\pm0.22$} & \textbf{31.54}{\scriptsize$\pm0.08$} & 22.25{\scriptsize$\pm0.45$} & 22.39{\scriptsize$\pm2.59$} \\
\midrule
\multicolumn{2}{l}{Completed seeds ($n/3$)} & 1/3 & 3/3 & 3/3 & 3/3 & 3/3 \\
\bottomrule
\end{tabular}
\caption{Six-arm ablation of the committed reference-model versions. Entries are completed-seed means and sample SD where at least two seeds are available. DTI uses the random held-out endpoint. Retrospective held-out means are multiplied by 100; completed seed counts are shown. Bold marks all displayed maxima. N/A denotes incomplete paired evaluation. Model versions were not selected by held-out scores; these endpoints are not blind confirmation.}
\label{tab:strong-ablation}
\end{table}

\begin{table}[t]
\centering\small
\setlength{\tabcolsep}{3pt}
\begin{tabular}{lrrrrr}
\toprule
Held-out contrast & DTI & Proteomics & VCC & Norman & Tahoe \\
\midrule
Loop $-$ fixed (1 lab) & N/A & -4.91 & -9.71 & -11.41 & -7.46 \\
Loop $-$ fixed (10 labs) & N/A & -0.08 & +0.00 & +6.81 & -1.49 \\
Loop $-$ direct (1 lab) & N/A & +0.00 & -0.87 & +0.00 & +0.00 \\
Loop $-$ direct (10 labs) & N/A & +0.00 & +0.00 & +0.00 & +0.00 \\
10 $-$ 1 labs (both fixed) & N/A & +12.02 & +1.61 & -10.07 & +5.97 \\
10 $-$ 1 labs (both direct) & N/A & +16.85 & +10.45 & +8.15 & +11.94 \\
10 $-$ 1 labs (both loop) & N/A & +16.85 & +11.32 & +8.15 & +11.94 \\
\midrule
Completed seeds ($n/3$) & 0/3 & 3/3 & 3/3 & 2/3 & 2/3 \\
\bottomrule
\end{tabular}
\caption{Separating total research/search improvement (loop minus fixed), feedback-specific improvement (loop minus matched-budget direct), and access to more laboratory data (ten minus one lab with loop fixed). Mean paired differences are in percentage points over the completed seeds on retrospective held-out data. Total search improvement must not be attributed solely to feedback. Negative effects and ties are retained; per-seed uncertainty is in Appendix R.}
\label{tab:strong-effects-full}
\end{table}

\begin{table}[t]
\centering\small
\setlength{\tabcolsep}{6pt}
\begin{tabular}{rrrrr}
\toprule
Proposal budget & Loop better & Tie & Direct better & Mean $\Delta$ primary \\
 & \multicolumn{3}{c}{task--seed trajectories} & points \\
\midrule
1 & 0 & 15 & 0 & +0.00 \\
2 & 2 & 13 & 0 & +0.78 \\
3 & 0 & 14 & 1 & -0.04 \\
4 & 0 & 15 & 0 & +0.00 \\
5 & 0 & 14 & 1 & -0.02 \\
6 & 0 & 15 & 0 & +0.00 \\
\bottomrule
\end{tabular}
\caption{Executed-prefix analysis of all 15 ten-laboratory development trajectories at common proposal budgets. At two proposals, feedback-guided search improves TAPB seed 44 by 4.84 points and Norman seed 42 by 6.90 points, with thirteen ties and no regressions.}
\label{tab:loop-prefix}
\end{table}

\clearpage
\begin{longtable}{llrrrr}
\caption{Held-out secondary metrics for the task-reference models, averaged over seeds 42--44 and multiplied by 100. Bold indicates all tied maxima, or minima for loss.}\label{tab:strong-secondary}\\
\toprule Method & Labs & Metric 1 & Metric 2 & Metric 3 & Metric 4 \\ \midrule\endfirsthead
\toprule Method & Labs & Metric 1 & Metric 2 & Metric 3 & Metric 4 \\ \midrule\endhead
\multicolumn{2}{l}{\textbf{DTI ($n=3$)}} & AUROC & AP & MCC & BCE \\
Fixed model & 1 & 82.83 & 81.57 & 49.95 & 52.37 \\
Direct & 1 & 89.89 & 87.77 & 65.71 & 61.02 \\
Our loop & 1 & 89.89 & 87.77 & 65.71 & 61.02 \\
Our loop & 10 & \textbf{93.76} & \textbf{92.73} & \textbf{73.71} & \textbf{36.59} \\
\midrule
\multicolumn{2}{l}{\textbf{Proteomics ($n=3$)}} & AP & AUROC & BCE & Acc. \\
Fixed model & 1 & 25.05 & 65.49 & 42.65 & \textbf{85.14} \\
Direct & 1 & 20.13 & 61.80 & 44.34 & 81.76 \\
Our loop & 1 & 20.13 & 61.80 & 44.34 & 81.76 \\
Our loop & 10 & \textbf{36.99} & \textbf{73.87} & \textbf{39.83} & \textbf{85.14} \\
\midrule
\multicolumn{2}{l}{\textbf{VCC ($n=3$)}} & Micro Top-1 & MRR & Top-3 & CE \\
Fixed model & 1 & 33.61 & \textbf{56.53} & \textbf{73.89} & \textbf{154.51} \\
Direct & 1 & 21.48 & 47.10 & 62.78 & 160.82 \\
Our loop & 1 & 20.74 & 46.94 & 64.26 & 160.86 \\
Our loop & 10 & \textbf{33.98} & 56.34 & 69.44 & 158.14 \\
\midrule
\multicolumn{2}{l}{\textbf{Norman ($n=3$)}} & Micro Top-1 & MRR & Top-3 & CE \\
Fixed model & 1 & \textbf{25.56} & 48.50 & 63.93 & \textbf{158.74} \\
Direct & 1 & 14.15 & 42.12 & 58.35 & 160.39 \\
Our loop & 1 & 14.15 & 42.12 & 58.35 & 160.39 \\
Our loop & 10 & 22.25 & \textbf{50.07} & \textbf{70.96} & 158.83 \\
\midrule
\multicolumn{2}{l}{\textbf{Tahoe ($n=3$)}} & Micro Top-1 & MRR & Top-3 & CE \\
Fixed model & 1 & 19.40 & 45.00 & 58.21 & 163.08 \\
Direct & 1 & 15.42 & 42.60 & 56.72 & 161.90 \\
Our loop & 1 & 15.42 & 42.60 & 56.72 & 161.90 \\
Our loop & 10 & \textbf{22.39} & \textbf{47.95} & \textbf{61.19} & \textbf{160.75} \\
\midrule
\bottomrule
\end{longtable}

\begin{longtable}{llrrl}
\caption{Per-seed paired cluster-bootstrap 95\% intervals for the committed model versions, in percentage points. These retrospective intervals condition on fixed trained models and do not measure across-seed uncertainty. DTI intervals, when available, are a post-hoc supplement using frozen predictions; random-endpoint rows use drug clusters.}\label{tab:strong-uncertainty}\\
\toprule Task & Contrast & Seed & Difference & 95\% interval \\ \midrule\endfirsthead
\toprule Task & Contrast & Seed & Difference & 95\% interval \\ \midrule\endhead
DTI & Loop $-$ fixed (10 labs) & 42 & +0.33 & [-0.06, +0.69] \\
DTI & Loop $-$ direct (1 lab) & 42 & +0.00 & [+0.00, +0.00] \\
DTI & Loop $-$ direct (10 labs) & 42 & +0.00 & [+0.00, +0.00] \\
DTI & 10 $-$ 1 labs (both fixed) & 42 & +12.51 & [+11.35, +13.71] \\
DTI & 10 $-$ 1 labs (both direct) & 42 & +5.78 & [+5.02, +6.57] \\
DTI & 10 $-$ 1 labs (both loop) & 42 & +5.78 & [+5.02, +6.57] \\
\midrule
DTI & Loop $-$ fixed (10 labs) & 43 & +0.54 & [-0.08, +1.11] \\
DTI & Loop $-$ direct (1 lab) & 43 & +0.00 & [+0.00, +0.00] \\
DTI & Loop $-$ direct (10 labs) & 43 & +0.00 & [+0.00, +0.00] \\
DTI & 10 $-$ 1 labs (both fixed) & 43 & +7.80 & [+6.82, +8.82] \\
DTI & 10 $-$ 1 labs (both direct) & 43 & +1.65 & [+0.90, +2.38] \\
DTI & 10 $-$ 1 labs (both loop) & 43 & +1.65 & [+0.90, +2.38] \\
\midrule
DTI & Loop $-$ fixed (10 labs) & 44 & +5.01 & [+4.25, +5.71] \\
DTI & Loop $-$ direct (1 lab) & 44 & +0.00 & [+0.00, +0.00] \\
DTI & Loop $-$ direct (10 labs) & 44 & +0.00 & [+0.00, +0.00] \\
DTI & 10 $-$ 1 labs (both fixed) & 44 & +6.60 & [+5.61, +7.74] \\
DTI & 10 $-$ 1 labs (both direct) & 44 & +4.17 & [+3.45, +4.93] \\
DTI & 10 $-$ 1 labs (both loop) & 44 & +4.17 & [+3.45, +4.93] \\
\midrule
Proteomics & Loop $-$ fixed (10 labs) & 42 & +0.00 & [+0.00, +0.00] \\
Proteomics & Loop $-$ direct (1 lab) & 42 & +0.00 & [+0.00, +0.00] \\
Proteomics & Loop $-$ direct (10 labs) & 42 & +0.00 & [+0.00, +0.00] \\
Proteomics & 10 $-$ 1 labs (both fixed) & 42 & +13.16 & [-0.39, +26.02] \\
Proteomics & 10 $-$ 1 labs (both direct) & 42 & +16.91 & [+3.37, +31.02] \\
Proteomics & 10 $-$ 1 labs (both loop) & 42 & +16.91 & [+3.37, +31.02] \\
\midrule
Proteomics & Loop $-$ fixed (10 labs) & 43 & +0.00 & [+0.00, +0.00] \\
Proteomics & Loop $-$ direct (1 lab) & 43 & +0.00 & [+0.00, +0.00] \\
Proteomics & Loop $-$ direct (10 labs) & 43 & +0.00 & [+0.00, +0.00] \\
Proteomics & 10 $-$ 1 labs (both fixed) & 43 & +15.98 & [+1.29, +29.00] \\
Proteomics & 10 $-$ 1 labs (both direct) & 43 & +17.03 & [+3.34, +31.44] \\
Proteomics & 10 $-$ 1 labs (both loop) & 43 & +17.03 & [+3.34, +31.44] \\
\midrule
Proteomics & Loop $-$ fixed (10 labs) & 44 & -0.24 & [-3.40, +2.33] \\
Proteomics & Loop $-$ direct (1 lab) & 44 & +0.00 & [+0.00, +0.00] \\
Proteomics & Loop $-$ direct (10 labs) & 44 & +0.00 & [+0.00, +0.00] \\
Proteomics & 10 $-$ 1 labs (both fixed) & 44 & +6.93 & [-3.43, +15.60] \\
Proteomics & 10 $-$ 1 labs (both direct) & 44 & +16.63 & [+3.02, +31.08] \\
Proteomics & 10 $-$ 1 labs (both loop) & 44 & +16.63 & [+3.02, +31.08] \\
\midrule
VCC & Loop $-$ fixed (10 labs) & 42 & +0.00 & [+0.00, +0.00] \\
VCC & Loop $-$ direct (1 lab) & 42 & +0.00 & [+0.00, +0.00] \\
VCC & Loop $-$ direct (10 labs) & 42 & +0.00 & [+0.00, +0.00] \\
VCC & 10 $-$ 1 labs (both fixed) & 42 & +1.70 & [-5.36, +8.58] \\
VCC & 10 $-$ 1 labs (both direct) & 42 & +9.90 & [-2.88, +23.04] \\
VCC & 10 $-$ 1 labs (both loop) & 42 & +9.90 & [-2.88, +23.04] \\
\midrule
VCC & Loop $-$ fixed (10 labs) & 43 & +0.00 & [+0.00, +0.00] \\
VCC & Loop $-$ direct (1 lab) & 43 & +4.83 & [-5.50, +16.49] \\
VCC & Loop $-$ direct (10 labs) & 43 & +0.00 & [+0.00, +0.00] \\
VCC & 10 $-$ 1 labs (both fixed) & 43 & +1.56 & [-5.76, +8.58] \\
VCC & 10 $-$ 1 labs (both direct) & 43 & +15.08 & [+2.94, +27.78] \\
VCC & 10 $-$ 1 labs (both loop) & 43 & +10.24 & [-2.74, +23.49] \\
\midrule
VCC & Loop $-$ fixed (10 labs) & 44 & +0.00 & [+0.00, +0.00] \\
VCC & Loop $-$ direct (1 lab) & 44 & -7.44 & [-17.23, +1.33] \\
VCC & Loop $-$ direct (10 labs) & 44 & +0.00 & [+0.00, +0.00] \\
VCC & 10 $-$ 1 labs (both fixed) & 44 & +1.56 & [-5.76, +8.58] \\
VCC & 10 $-$ 1 labs (both direct) & 44 & +6.38 & [-4.38, +19.13] \\
VCC & 10 $-$ 1 labs (both loop) & 44 & +13.82 & [+2.54, +26.34] \\
\midrule
Norman & Loop $-$ fixed (10 labs) & 42 & +7.26 & [-9.41, +25.41] \\
Norman & Loop $-$ direct (1 lab) & 42 & +0.00 & [+0.00, +0.00] \\
Norman & Loop $-$ direct (10 labs) & 42 & +0.00 & [+0.00, +0.00] \\
Norman & 10 $-$ 1 labs (both fixed) & 42 & -10.07 & [-21.19, +0.00] \\
Norman & 10 $-$ 1 labs (both direct) & 42 & +8.59 & [-8.30, +26.00] \\
Norman & 10 $-$ 1 labs (both loop) & 42 & +8.59 & [-8.30, +26.00] \\
\midrule
Norman & Loop $-$ fixed (10 labs) & 43 & +6.37 & [-11.12, +24.81] \\
Norman & Loop $-$ direct (1 lab) & 43 & +0.00 & [+0.00, +0.00] \\
Norman & Loop $-$ direct (10 labs) & 43 & +0.00 & [+0.00, +0.00] \\
Norman & 10 $-$ 1 labs (both fixed) & 43 & -10.07 & [-21.19, +0.00] \\
Norman & 10 $-$ 1 labs (both direct) & 43 & +7.70 & [-9.33, +25.85] \\
Norman & 10 $-$ 1 labs (both loop) & 43 & +7.70 & [-9.33, +25.85] \\
\midrule
Norman & Loop $-$ fixed (10 labs) & 44 & +6.67 & [-10.37, +24.81] \\
Norman & Loop $-$ direct (1 lab) & 44 & +0.00 & [+0.00, +0.00] \\
Norman & Loop $-$ direct (10 labs) & 44 & +0.00 & [+0.00, +0.00] \\
Norman & 10 $-$ 1 labs (both fixed) & 44 & -10.07 & [-21.19, +0.00] \\
Norman & 10 $-$ 1 labs (both direct) & 44 & +8.00 & [-8.52, +25.93] \\
Norman & 10 $-$ 1 labs (both loop) & 44 & +8.00 & [-8.52, +25.93] \\
\midrule
Tahoe & Loop $-$ fixed (10 labs) & 42 & +0.00 & [+0.00, +0.00] \\
Tahoe & Loop $-$ direct (1 lab) & 42 & +0.00 & [+0.00, +0.00] \\
Tahoe & Loop $-$ direct (10 labs) & 42 & +0.00 & [+0.00, +0.00] \\
Tahoe & 10 $-$ 1 labs (both fixed) & 42 & +4.48 & [-5.97, +14.93] \\
Tahoe & 10 $-$ 1 labs (both direct) & 42 & +10.45 & [-1.49, +22.39] \\
Tahoe & 10 $-$ 1 labs (both loop) & 42 & +10.45 & [-1.49, +22.39] \\
\midrule
Tahoe & Loop $-$ fixed (10 labs) & 43 & -2.99 & [-7.46, +0.00] \\
Tahoe & Loop $-$ direct (1 lab) & 43 & +0.00 & [+0.00, +0.00] \\
Tahoe & Loop $-$ direct (10 labs) & 43 & +0.00 & [+0.00, +0.00] \\
Tahoe & 10 $-$ 1 labs (both fixed) & 43 & +7.46 & [-4.48, +17.95] \\
Tahoe & 10 $-$ 1 labs (both direct) & 43 & +13.43 & [+1.49, +23.88] \\
Tahoe & 10 $-$ 1 labs (both loop) & 43 & +13.43 & [+1.49, +23.88] \\
\midrule
Tahoe & Loop $-$ fixed (10 labs) & 44 & -7.46 & [-19.40, +4.48] \\
Tahoe & Loop $-$ direct (1 lab) & 44 & +0.00 & [+0.00, +0.00] \\
Tahoe & Loop $-$ direct (10 labs) & 44 & +0.00 & [+0.00, +0.00] \\
Tahoe & 10 $-$ 1 labs (both fixed) & 44 & +7.46 & [-4.48, +17.95] \\
Tahoe & 10 $-$ 1 labs (both direct) & 44 & -2.99 & [-17.91, +11.94] \\
Tahoe & 10 $-$ 1 labs (both loop) & 44 & -2.99 & [-17.91, +11.94] \\
\midrule
\bottomrule
\end{longtable}

\begin{longtable}{llrrr}
\caption{All three frozen TAPB held-out endpoints under the committed common protocol. Scores are completed-seed means multiplied by 100; N/A means no complete independently verified six-arm evaluation. The main DTI endpoint is random held-out.}\label{tab:strong-dti-endpoints}\\
\toprule Endpoint & Method & Labs & AUROC & AP \\ \midrule\endfirsthead
\toprule Endpoint & Method & Labs & AUROC & AP \\ \midrule\endhead
Random & Fixed model & 1 & 81.76 & 82.30 \\
Random & Direct & 1 & 88.82 & 85.98 \\
Random & Our loop & 1 & 88.82 & 85.98 \\
Random & Fixed model & 10 & 94.27 & 92.87 \\
Random & Direct & 10 & \textbf{94.60} & \textbf{94.27} \\
Random & Our loop & 10 & \textbf{94.60} & \textbf{94.27} \\
\midrule
Unseen drug & Fixed model & 1 & 80.67 & 86.41 \\
Unseen drug & Direct & 1 & 86.68 & 87.08 \\
Unseen drug & Our loop & 1 & 86.68 & 87.08 \\
Unseen drug & Fixed model & 10 & 92.96 & 93.58 \\
Unseen drug & Direct & 10 & \textbf{93.06} & \textbf{94.63} \\
Unseen drug & Our loop & 10 & \textbf{93.06} & \textbf{94.63} \\
\midrule
Unseen protein & Fixed model & 1 & 69.86 & 78.05 \\
Unseen protein & Direct & 1 & 81.05 & 77.04 \\
Unseen protein & Our loop & 1 & 81.05 & 77.04 \\
Unseen protein & Fixed model & 10 & 89.88 & 86.19 \\
Unseen protein & Direct & 10 & \textbf{89.91} & \textbf{90.30} \\
Unseen protein & Our loop & 10 & \textbf{89.91} & \textbf{90.30} \\
\midrule
\bottomrule
\end{longtable}

\begingroup\scriptsize\setlength{\tabcolsep}{3pt}
\begin{longtable}{llrrrr}
\caption{Post-hoc TAPB uncertainty supplement from frozen prediction files. All three endpoints, six arms and fixed contrasts are retained. Entries are paired AUROC/AP differences and percentile 95\% intervals in percentage points from 1,000 requested cluster draws (RNG seed 20260918). Random and unseen-drug endpoints resample drug clusters; unseen-protein resamples protein clusters. Intervals condition on fixed trained models and one clustering axis; they are not two-way, across-seed, multiplicity-adjusted, or prospective confirmatory intervals.}\label{tab:strong-dti-uncertainty}\\
\toprule Endpoint / seed & Contrast & AUROC $\Delta$ & 95\% interval & AP $\Delta$ & 95\% interval \\ \midrule\endfirsthead
\toprule Endpoint / seed & Contrast & AUROC $\Delta$ & 95\% interval & AP $\Delta$ & 95\% interval \\ \midrule\endhead
Random / 42 & Loop $-$ fixed (10 labs) & +0.33 & [-0.06, +0.69] & +1.40 & [+0.76, +2.05] \\
Random / 42 & Loop $-$ direct (1 lab) & +0.00 & [+0.00, +0.00] & +0.00 & [+0.00, +0.00] \\
Random / 42 & Loop $-$ direct (10 labs) & +0.00 & [+0.00, +0.00] & +0.00 & [+0.00, +0.00] \\
Random / 42 & 10 $-$ 1 labs (fixed) & +12.51 & [+11.35, +13.71] & +10.56 & [+9.10, +12.11] \\
Random / 42 & 10 $-$ 1 labs (direct) & +5.78 & [+5.02, +6.57] & +8.29 & [+7.01, +9.56] \\
Random / 42 & 10 $-$ 1 labs (loop) & +5.78 & [+5.02, +6.57] & +8.29 & [+7.01, +9.56] \\
Random / 42 & Our harness $-$ fixed (1 lab) & +12.84 & [+11.82, +14.01] & +11.96 & [+10.56, +13.45] \\
Random / 42 & Our harness $-$ direct (1 lab) & +5.78 & [+5.02, +6.57] & +8.29 & [+7.01, +9.56] \\
\midrule
Unseen drug / 42 & Loop $-$ fixed (10 labs) & +0.10 & [-0.47, +0.67] & +1.06 & [+0.42, +1.66] \\
Unseen drug / 42 & Loop $-$ direct (1 lab) & +0.00 & [+0.00, +0.00] & +0.00 & [+0.00, +0.00] \\
Unseen drug / 42 & Loop $-$ direct (10 labs) & +0.00 & [+0.00, +0.00] & +0.00 & [+0.00, +0.00] \\
Unseen drug / 42 & 10 $-$ 1 labs (fixed) & +12.30 & [+10.77, +13.98] & +7.17 & [+5.82, +8.72] \\
Unseen drug / 42 & 10 $-$ 1 labs (direct) & +6.38 & [+5.34, +7.53] & +7.55 & [+6.30, +8.92] \\
Unseen drug / 42 & 10 $-$ 1 labs (loop) & +6.38 & [+5.34, +7.53] & +7.55 & [+6.30, +8.92] \\
Unseen drug / 42 & Our harness $-$ fixed (1 lab) & +12.40 & [+10.94, +13.85] & +8.22 & [+6.97, +9.66] \\
Unseen drug / 42 & Our harness $-$ direct (1 lab) & +6.38 & [+5.34, +7.53] & +7.55 & [+6.30, +8.92] \\
\midrule
Unseen protein / 42 & Loop $-$ fixed (10 labs) & +0.03 & [-5.53, +4.28] & +4.12 & [-4.10, +8.22] \\
Unseen protein / 42 & Loop $-$ direct (1 lab) & +0.00 & [+0.00, +0.00] & +0.00 & [+0.00, +0.00] \\
Unseen protein / 42 & Loop $-$ direct (10 labs) & +0.00 & [+0.00, +0.00] & +0.00 & [+0.00, +0.00] \\
Unseen protein / 42 & 10 $-$ 1 labs (fixed) & +20.02 & [+0.16, +63.83] & +8.14 & [-8.27, +51.77] \\
Unseen protein / 42 & 10 $-$ 1 labs (direct) & +8.86 & [+2.48, +13.39] & +13.26 & [+3.15, +28.50] \\
Unseen protein / 42 & 10 $-$ 1 labs (loop) & +8.86 & [+2.48, +13.39] & +13.26 & [+3.15, +28.50] \\
Unseen protein / 42 & Our harness $-$ fixed (1 lab) & +20.05 & [+1.84, +61.27] & +12.25 & [-3.01, +53.19] \\
Unseen protein / 42 & Our harness $-$ direct (1 lab) & +8.86 & [+2.48, +13.39] & +13.26 & [+3.15, +28.50] \\
\midrule
Random / 43 & Loop $-$ fixed (10 labs) & +0.54 & [-0.08, +1.11] & +0.44 & [-0.54, +1.44] \\
Random / 43 & Loop $-$ direct (1 lab) & +0.00 & [+0.00, +0.00] & +0.00 & [+0.00, +0.00] \\
Random / 43 & Loop $-$ direct (10 labs) & +0.00 & [+0.00, +0.00] & +0.00 & [+0.00, +0.00] \\
Random / 43 & 10 $-$ 1 labs (fixed) & +7.80 & [+6.82, +8.82] & +7.68 & [+6.32, +9.13] \\
Random / 43 & 10 $-$ 1 labs (direct) & +1.65 & [+0.90, +2.38] & +1.17 & [-0.11, +2.50] \\
Random / 43 & 10 $-$ 1 labs (loop) & +1.65 & [+0.90, +2.38] & +1.17 & [-0.11, +2.50] \\
Random / 43 & Our harness $-$ fixed (1 lab) & +8.35 & [+7.29, +9.37] & +8.12 & [+6.53, +9.76] \\
Random / 43 & Our harness $-$ direct (1 lab) & +1.65 & [+0.90, +2.38] & +1.17 & [-0.11, +2.50] \\
\midrule
Unseen drug / 43 & Loop $-$ fixed (10 labs) & -0.51 & [-1.44, +0.29] & -1.10 & [-2.16, -0.18] \\
Unseen drug / 43 & Loop $-$ direct (1 lab) & +0.00 & [+0.00, +0.00] & +0.00 & [+0.00, +0.00] \\
Unseen drug / 43 & Loop $-$ direct (10 labs) & +0.00 & [+0.00, +0.00] & +0.00 & [+0.00, +0.00] \\
Unseen drug / 43 & 10 $-$ 1 labs (fixed) & +6.68 & [+5.33, +8.25] & +4.66 & [+3.33, +6.36] \\
Unseen drug / 43 & 10 $-$ 1 labs (direct) & -0.15 & [-0.96, +0.67] & -1.43 & [-2.43, -0.42] \\
Unseen drug / 43 & 10 $-$ 1 labs (loop) & -0.15 & [-0.96, +0.67] & -1.43 & [-2.43, -0.42] \\
Unseen drug / 43 & Our harness $-$ fixed (1 lab) & +6.17 & [+4.85, +7.48] & +3.57 & [+2.10, +5.09] \\
Unseen drug / 43 & Our harness $-$ direct (1 lab) & -0.15 & [-0.96, +0.67] & -1.43 & [-2.43, -0.42] \\
\midrule
Unseen protein / 43 & Loop $-$ fixed (10 labs) & -0.96 & [-7.24, +3.07] & -1.83 & [-21.55, +7.54] \\
Unseen protein / 43 & Loop $-$ direct (1 lab) & +0.00 & [+0.00, +0.00] & +0.00 & [+0.00, +0.00] \\
Unseen protein / 43 & Loop $-$ direct (10 labs) & +0.00 & [+0.00, +0.00] & +0.00 & [+0.00, +0.00] \\
Unseen protein / 43 & 10 $-$ 1 labs (fixed) & +4.79 & [-6.61, +28.12] & +0.41 & [-16.48, +31.69] \\
Unseen protein / 43 & 10 $-$ 1 labs (direct) & -1.19 & [-6.22, +6.09] & -6.76 & [-13.95, +6.71] \\
Unseen protein / 43 & 10 $-$ 1 labs (loop) & -1.19 & [-6.22, +6.09] & -6.76 & [-13.95, +6.71] \\
Unseen protein / 43 & Our harness $-$ fixed (1 lab) & +3.83 & [-5.83, +22.32] & -1.42 & [-14.49, +16.88] \\
Unseen protein / 43 & Our harness $-$ direct (1 lab) & -1.19 & [-6.22, +6.09] & -6.76 & [-13.95, +6.71] \\
\midrule
\bottomrule
\end{longtable}
\endgroup


\paragraph{Uncertainty analysis.}
The DTI analysis is computed after training and held-out scoring, using paired cluster bootstrap with 1,000 fixed-seed resamples of retained predictions. It groups by released drug identity on random/unseen-drug endpoints and by protein sequence on unseen-protein. Other task intervals use their natural compound or intervention grouping. Resamples are shared within each method contrast. These intervals characterize sample-level variation within each selected fit; between-training-seed variation is reported through the three-seed means and sample SD.

\paragraph{Verification.}
CPU unit and integration tests cover the common trainer, task adapters, proposal generation, frozen evaluation and cell attention-mask correction. Publication checks bind table cells to completed arms, model versions, checkpoint seals and independent score recomputation. The snapshot includes receipt hashes, aggregate metrics and the complete proposal histories used in this manuscript.
